% !TeX program = xelatex
\documentclass[11pt,a4paper]{article}
\usepackage{fontspec}
\setmainfont{DejaVu Serif}
\setsansfont{DejaVu Sans}
\setmonofont{DejaVu Sans Mono}
% Polski bez babel/polyglossia (brak texlive-lang-polish w środowisku kompilacji):
\renewcommand{\contentsname}{Spis treści}
\renewcommand{\tablename}{Tabela}
\renewcommand{\figurename}{Rysunek}
\hyphenpenalty=3000 \tolerance=2000 \emergencystretch=2em
\usepackage[margin=2.1cm]{geometry}
\usepackage{amsmath}
\usepackage{xcolor}
\definecolor{granat}{RGB}{12,35,88}
\usepackage{booktabs}
\usepackage{tabularx}
\usepackage{longtable}
\usepackage{array}
\usepackage{enumitem}
\usepackage[colorlinks=true,linkcolor=granat,urlcolor=granat]{hyperref}
\usepackage{titlesec}
\usepackage{parskip}
\titleformat{\section}{\large\bfseries\sffamily\color{granat}}{\thesection.}{0.6em}{}
\titleformat{\subsection}{\normalsize\bfseries\sffamily\color{granat}}{\thesubsection.}{0.6em}{}
\titleformat{\subsubsection}{\normalsize\bfseries\sffamily\color{granat!80!black}}{\thesubsubsection.}{0.6em}{}
\setlist{nosep,leftmargin=1.4em}
\newcolumntype{L}{>{\raggedright\arraybackslash}X}
\newcommand{\ohm}{$\Omega$}
\newcommand{\popr}[1]{\textbf{\color{granat}[POPRAWKA v3 --- #1]}}

\title{\textbf{\color{granat}AMPERE POINT --- custom device}\\[4pt]
\large Objaśnienie schematów elektrycznych (arkusze E1--E3)\\
z przeglądem poprawności i wykazem poprawek wprowadzonych w wersji 3\\[2pt]
\normalsize wersja 1}
\author{Opracowanie wewnętrzne AMPERE POINT}
\date{6 lipca 2026}

\begin{document}
\maketitle
\vspace{-1.6em}

\noindent\textbf{Status dokumentu:} objaśnienie techniczne do schematu elektrycznego custom device
(\texttt{schemat\_elektryczny\_v1}, \texttt{v2\_EPLAN} z 2026-07-03 oraz \texttt{v3\_EPLAN} z 2026-07-06
--- wersja 3 powstała razem z niniejszym dokumentem i wprowadza poprawki opisane w rozdz.~\ref{sec:poprawki}).
Dokument pisany dla osoby z podstawami elektrotechniki; wyjaśnia \emph{każdy} aparat na schemacie:
co to jest, po co jest, jak działa i jak jest podłączony --- ze szczególnym naciskiem na wszystkie
urządzenia dołączone do mikrokontrolera (rozdz.~\ref{sec:mcu}).

\tableofcontents

% =====================================================================
\section{Czym jest ten schemat i skąd się wziął}
\label{sec:wstep}

Custom device to zautomatyzowana stacja testowa ładowarek AC (EVSE), która wobec badanej ładowarki
\textbf{udaje samochód elektryczny}: symuluje stany złącza wg IEC~61851-1, pobiera realny prąd
testowy (do 32~A na fazę) i rejestruje wyniki. Umożliwia wykonanie pięciu obowiązkowych pomiarów
z ,,Przewodnika w zakresie wykonywania pomiarów elektrycznych stacji ładowania'' (Warszawa 2024)
oraz prób funkcjonalnych pod obciążeniem --- bez użycia auta testowego.

Schemat elektryczny jest bezpośrednim przełożeniem zatwierdzonej koncepcji
(\texttt{AMPERE\_POINT\_custom\_device\_koncepcja\_v1.pdf}, punkty Z1--Z6) i schematów ideowych v1.
Kluczowe decyzje, które ukształtowały topologię:

\begin{itemize}
\item \textbf{Z1 --- przelot mocy z odczepem sterującym.} Prąd obciążenia \emph{nie płynie} przez
tester PM701E (adaptery tej klasy nie są projektowane na ciągłe 32~A; porównywalny EVSE-200
dopuszcza $\sim$10~A). Urządzenie ma własne gniazdo pojazdowe Type~2 i własny tor mocy 6~mm$^2$,
a PM701E wisi na krótkim odczepie i steruje wyłącznie sygnałami.
\item \textbf{Z2 --- architektura dwumodułowa.} Moduł~A (sterowanie + pomiary, rozdzielnica IP54)
i moduł~B (obciążenie --- grzałki z wentylatorami) połączone przewodem 5$\times$6~mm$^2$ z wtykiem
CEE 32~A~5p oraz osobnym przewodem sterowniczym (złącze -X5).
\item \textbf{Bezpieczeństwo sprzętowe, nie programowe.} Załączenie mocy warunkuje szeregowy
łańcuch styków 24~V (arkusz E2), którego firmware nie może zmostkować --- mikrokontroler może
jedynie \emph{zezwalać} i \emph{rozłączać}.
\end{itemize}

Schemat składa się z trzech arkuszy, które odpowiadają trzem warstwom urządzenia:
\textbf{E1} --- tor mocy (prąd testowy), \textbf{E2} --- łańcuch bezpieczeństwa (zezwolenie na moc),
\textbf{E3} --- sterowanie, pomiary i moduł obciążenia (mózg i mięśnie).

% =====================================================================
\section{Jak czytać schemat --- konwencje}
\label{sec:konwencje}

Wersja 2/3 schematu odwzorowuje konwencję dokumentacji EPLAN (symbole IEC~60617):

\begin{itemize}
\item \textbf{Rama arkusza} ma siatkę odniesień: kolumny 1--10 i wiersze A--F. Pozycję aparatu
opisuje się jako \emph{/arkusz.kolumna}, np. \texttt{/E2.3} = arkusz E2, kolumna 3.
\item \textbf{Odsyłacze krzyżowe}: przy styku podany jest adres cewki, która nim porusza; pod cewką
znajduje się \emph{lustro styków} --- spis wszystkich styków danego aparatu z adresami, gdzie
zostały narysowane. Dzięki temu aparat rozrysowany na kilku arkuszach (np. stycznik -K1: cewka na
E2, styki główne na E1) pozostaje jednoznaczny.
\item \textbf{Oznaczenia referencyjne} (litera wg IEC~81346): \texttt{-X}~złącza/zaciski,
\texttt{-Q}~łączniki mocy, \texttt{-F}~zabezpieczenia, \texttt{-K}~styczniki/przekaźniki,
\texttt{-B}~czujniki, \texttt{-T}~przekładniki, \texttt{-E}~grzałki, \texttt{-M}~silniki/serwa/wentylatory,
\texttt{-A}~zespoły (płyta logiki, moduły), \texttt{-G}~zasilacze, \texttt{-P}~HMI/wskaźniki,
\texttt{-S}~łączniki sterownicze, \texttt{-U}~układy scalone/przetworniki, \texttt{-Y}~elektromagnesy.
\item \textbf{Numeracja zacisków}: tory główne łączników 1-2, 3-4, 5-6, 7-8; styki pomocnicze
zwierne x3-x4 (np. 13-14, 23-24), rozwierne x1-x2 (np. 11-12, 21-22); cewka A1-A2.
\emph{Uwaga: zaciski cewki A1/A2 to co innego niż oznaczenia urządzeń -A1 (płyta ESP32) i -A2
(watchdog) --- zbieżność jest przypadkowa i wynika z normy.}
\item \textbf{Stany aparatów} rysowane są zawsze w stanie beznapięciowym/spoczynkowym
(cewki niewzbudzone, przyciski niewciśnięte).
\end{itemize}

% =====================================================================
\section{Arkusz E1 --- tor mocy 3F}
\label{sec:e1}

Tor mocy prowadzi energię z badanej ładowarki (DUT) do modułu obciążenia:
\[
\text{-X1 (Type 2)} \;\to\; \text{-Q0} \;\to\; \text{-F1..3} \;\to\; \text{-B1} \;\to\;
\text{-T1..3} \;\to\; \text{-K1} \;\to\; \text{-X2 (CEE)} \;\to\; \text{moduł B}
\]
z odczepem \textbf{-X3} do PM701E pobranym bezpośrednio przy gnieździe -X1.

\subsection{-X1 --- gniazdo pojazdowe Type 2 32 A 3F}
\textbf{Funkcja:} wejście urządzenia --- tu wpina się wtyk kabla badanej ładowarki, dokładnie tak,
jak wpinałby się w gniazdo samochodu. \textbf{Styki:} L1, L2, L3, N, PE (tor mocy) oraz CP i PP
(sygnałowe). \textbf{Działanie:} gniazdo jest stroną ,,pojazdową'' złącza wg IEC~62196 --- dzięki
temu z punktu widzenia ładowarki urządzenie jest elektrycznie nieodróżnialne od auta. Wersja
z \textbf{blokadą elektromagnesową -Y1} (rozdz.~\ref{sec:y1}) uniemożliwia wyciągnięcie wtyku
w trakcie testu pod napięciem. Dla ładowarek z gniazdem (zamiast kabla) używa się przewodu
Type~2--Type~2; dla Type~1 --- posiadanej przejściówki.

\subsection{-X3 --- odczep sterujący do PM701E}
\label{sec:x3}
\textbf{Funkcja:} przenosi sygnały i zasilanie własne do testera PM701E, który jest ,,symulatorem
stanów pojazdu'' całego układu. \textbf{Działanie:} PM701E swoimi pokrętłami (stan BB/C/A/B/D,
prąd NC/13/20/32/63~A) ustawia rezystancje na linii CP zgodnie z IEC~61851-1 --- ładowarka
,,widzi'' samochód w zadanym stanie. Pokrętłami fizycznie obracają serwa -M1/-M2 (rozdz.~\ref{sec:serwa}).
Gniazda bananowe PM701E (L1/N/PE/L2/L3) pozostają dostępne z zewnątrz jako punkt przyłączenia
miernika UT595 do pomiarów protokołowych --- elektrycznie to ten sam węzeł co magistrala mocy,
bo odczep łączy się z nią wprost na zaciskach -X1. Para SIGNAL~OUTPUT (CP + masa) biegnie do toru
akwizycji CP mikrokontrolera (rozdz.~\ref{sec:cp}).

\textbf{Pobór prądu:} przez odczep płynie tylko pobór własny PM701E (pomijalny) i prądy pomiarowe
przyrządów --- \emph{nie} prąd obciążenia. Dlatego wystarcza przewód 1,5~mm$^2$ o długości
$\le$1~m.

\popr{odczep 7-żyłowy} W v1/v2 odczep miał 5 żył (L1, N, PE, CP, PP) --- gniazda bananowe L2/L3
PM701E byłyby martwe, wbrew koncepcji (rozdz.~B1: ,,przewód 5$\times$1,5~mm$^2$ + CP/PP'', czyli
7 żył) i wbrew sensowi pomiarów 3-fazowych przy Q11. W v3 odczep prowadzi
\textbf{L1, L2, L3, N, PE + CP, PP}.

\subsection{-Q0 --- rozłącznik izolacyjny 4P 40 A}
\textbf{Funkcja:} ręczne, widoczne odcięcie toru mocy do prac przyłączeniowych (wymiana modułu~B,
prace na CEE). \textbf{Działanie:} rozłącznik (nie wyłącznik) --- służy do przerywania obwodu,
rozłącza wszystkie 4 bieguny (L1/L2/L3/N) jednocześnie; nie ma zdolności wyłączania zwarć (od tego
są -F1..3). \textbf{Ważne:} -Q0 leży \emph{za} odczepem -X3, więc \textbf{nie odłącza PM701E ani
gniazd bananowych} --- przy wpiętym DUT w stanie~C pozostają one pod napięciem niezależnie od
położenia -Q0 i stanu -K1. To cecha topologii (gniazda bananowe muszą widzieć napięcie do pomiarów
Zs/RCD przy otwartym -K1), ale wymaga świadomości operatora. \popr{adnotacja} W v3 dodano
adnotację ostrzegawczą na arkuszu i wymóg tabliczki przy gniazdach bananowych; pełną,
sprzętową odpowiedź na ten problem daje też nowy przekaźnik -K4 (rozdz.~\ref{sec:k4}).

\subsection{-F1, -F2, -F3 --- bezpieczniki topikowe gG 32 A}
\textbf{Funkcja:} zabezpieczenie zwarciowe i przeciążeniowe toru mocy urządzenia.
\textbf{Działanie:} wkładka topikowa gG (pełnozakresowa, ogólnego przeznaczenia) przepala się
przy przeciążeniu/zwarciu. Po jednej na fazę; \textbf{przewód N celowo bez zabezpieczenia} ---
przerwanie samego N pod obciążeniem 3F wywołałoby niebezpieczne przepięcia asymetryczne;
N jest za to rozłączany razem z fazami (czterobiegunowe -Q0 i -K1). PE nie wolno przerywać niczym
--- biegnie nieprzerwanie od -X1 do -X2.

\subsection{-B1 --- czujnik prądu różnicowego AC/DC (CDSR, funkcja RCM)}
\label{sec:b1}
\textbf{Funkcja:} ciągły monitor prądu różnicowego (Residual Current Monitor) własnych torów
urządzenia --- ochrona operatora i sprzętu; dodatkowo weryfikacja krzyżowa generatora upływu -A3.
\textbf{Zasada działania:} przez okno czujnika przechodzą \emph{wszystkie} przewody robocze
(L1+L2+L3+N --- \textbf{bez PE!}). W obwodzie sprawnym suma chwilowych prądów = 0; każda różnica
oznacza, że część prądu wraca poza torem (np. przez PE lub człowieka). Zwykły przekładnik widzi
tylko różnicę przemienną; CDSR to czujnik \textbf{fluxgate} (bramka strumieniowa): rdzeń jest
cyklicznie wprowadzany w nasycenie przebiegiem pomocniczym, a składowa stała prądu różnicowego
przesuwa symetrię nasycania --- dzięki temu czujnik widzi także \textbf{gładki prąd stały}
(technologia z prawdziwych RCD typu~B / modułów RDC-DD). Typ: LEM \textbf{CDSR~0.07-TP},
zakres $\pm$70~mA, rozdzielczość $\sim$0,5~mA, zasilanie 3,3~V, wyjście \textbf{SPI} (+ analogowe)
wprost do -A1 (rozdz.~\ref{sec:mcu-spi}).
\textbf{Reakcja:} przekroczenie progu (np. 20~mA AC / 5~mA DC nieoczekiwanego pochodzenia)
$\to$ firmware zdejmuje zezwolenie -K10 $\to$ rozpad łańcucha $\to$ -K1 otwiera. Podczas testu
RDC-DD (celowe wymuszanie upływu przez -A3) firmware podnosi/maskuje próg i jednocześnie
\emph{porównuje} odczyt CDSR z prądem zadanym --- to weryfikacja krzyżowa obu bloków.

\subsection{-T1, -T2, -T3 --- przekładniki prądowe 40 A / 1 V}
\textbf{Funkcja:} pomiar prądów fazowych pobieranych przez moduł obciążenia.
\textbf{Działanie:} klasyczne przekładniki przelotowe z wyjściem napięciowym (wbudowany rezystor
obciążający): 40~A po stronie pierwotnej $\to$ 1~V na wyjściu. Wyjścia idą do wejść prądowych
układu pomiaru energii -U2 (ATM90E36A, rozdz.~\ref{sec:u2}). Uwaga: CT mierzy tylko AC --- to
wystarcza, bo prąd obciążenia jest sinusoidalny; składową DC upływu mierzy -B1 i bocznik -A3.

\subsection{-K1 --- stycznik główny 4P 40 A (AC-1)}
\textbf{Funkcja:} jedyny aparat, który podaje moc do modułu obciążenia; wykonawczy element
łańcucha bezpieczeństwa. \textbf{Działanie:} cztery styki główne (L1/L2/L3/N) zamykane
elektromagnesem; cewka 24~V zasilana \emph{wyłącznie} przez styk przekaźnika bezpieczeństwa -RB
(arkusz E2, rozdz.~\ref{sec:e2}). Kategoria \textbf{AC-1} (obciążenia rezystancyjne) jest właściwa,
bo grzałki nie mają udarów rozruchowych ani indukcyjności; 40~A daje zapas względem 32~A.
Rozpad łańcucha = otwarcie -K1 = moduł~B beznapięciowy. Magistrala \emph{przed} -K1 (gniazda
bananowe, -B1, -T) pozostaje pod napięciem, jeśli DUT jest w stanie~C --- patrz -Q0 i -K4.

\subsection{-X2 --- gniazdo CEE 32 A 5p (interfejs modułu B)}
\textbf{Funkcja:} wyjście toru mocy do modułu obciążenia. Przewód gumowy \textbf{5$\times$6~mm$^2$
H07RN-F} z wtykiem CEE 32~A~5p --- standard przemysłowy, jednoznaczne przyporządkowanie
L1/L2/L3/N/PE, odporność mechaniczna. Przekrój 6~mm$^2$ z zapasem dla 32~A pracy ciągłej.

\subsection{Nowe elementy E1 w wersji 3}
\label{sec:k4}
\popr{-K4, przerwanie CP} \textbf{Problem:} E-STOP w v1/v2 otwierał tylko -K1 --- odcinał moduł~B,
ale magistrala i gniazda bananowe pozostawały pod napięciem od DUT będącego w stanie~C.
\textbf{Rozwiązanie:} przekaźnik \textbf{-K4} z zestykiem w linii \textbf{CP odczepu -X3}.
Cewka -K4 jest zasilana z 24~V przez \emph{drugi} zestyk rozwierny E-STOPu (-S0:21-22).
Wciśnięcie E-STOP (albo zanik 24~V) $\to$ -K4 odpada $\to$ linia CP zostaje przerwana $\to$
ładowarka widzi odłączenie pojazdu (stan~A) i \textbf{sama otwiera swój stycznik} --- magistrala
L1/L2/L3/N robi się martwa w czasie wymaganym przez IEC~61851 (typ.\ $\le$100~ms dla utraty CP).
To ten sam mechanizm, którego używa przycisk ,,CP~Error~E'' na PM701E, tylko wpięty na stałe
w obwód bezpieczeństwa. Konsekwencja eksploatacyjna: przy niezasilonym urządzeniu nie da się
osiągnąć stanu~C (CP przerwane) --- zachowanie celowe i bezpieczne; na listwie -X3 przewidziano
mostek serwisowy do świadomego obejścia (np. użycie samego PM701E bez custom device).

\popr{punkt przyłączenia -A3} Doprecyzowano miejsce odczepu generatora upływu: \textbf{L1 za oknem
-B1} (między -B1 a -T1). Powód: prąd upływu wymuszany przez -A3 musi przechodzić przez okno -B1,
inaczej niemożliwa byłaby weryfikacja krzyżowa generatora czujnikiem CDSR (wymóg koncepcji,
rozdz.~B5b) i nadzór RCM nad testem. Dla samego DUT miejsce odczepu za -B1 niczego nie zmienia
(upływ i tak wraca przez PE poza czujnikiem RDC-DD ładowarki, więc ładowarka go widzi).

% =====================================================================
\section{Arkusz E2 --- łańcuch bezpieczeństwa 24 V}
\label{sec:e2}

\subsection{Zasada: obwód prądu spoczynkowego}
Cewka stycznika głównego -K1 może dostać 24~V \textbf{wyłącznie} przez szereg styków fizycznych
aparatów. To obwód \emph{prądu spoczynkowego} (fail-safe): przerwa w dowolnym miejscu ---
zadziałanie zabezpieczenia, awaria aparatu, urwany przewód, zanik 24~V, zawieszenie firmware ---
zawsze prowadzi do stanu bezpiecznego (rozpad -RB $\to$ otwarcie -K1). Nie istnieje pojedyncza
usterka, która ,,sama z siebie'' załączy moc; firmware nie jest w stanie zmostkować żadnego członu,
bo jego jedyny wpływ na łańcuch to styki -K10 i -A2 (oba mogą obwód tylko \emph{rozewrzeć}).

\subsection{Człony łańcucha (szczebel 1: zezwolenie $\to$ cewka -RB)}
Kolejno od szyny +24~V:
\begin{enumerate}
\item \textbf{-S0 E-STOP} --- grzybek dłoniowy, zestyk rozwierny (NC) z ryglowaniem; na panelu
modułu~A. Wciśnięcie = natychmiastowa przerwa. Od v3 drugi zestyk -S0:21-22 zwalnia też -K4
(przerwanie CP, rozdz.~\ref{sec:k4}), więc E-STOP gasi \emph{całą} magistralę, nie tylko moduł~B.
\item \textbf{-F11.1 \ldots -F11.12 --- termiki KSD301} --- bimetalowe wyłączniki termiczne
($\sim$90~$^\circ$C, NC) przykręcone do każdej sekcji grzałek w module~B, połączone szeregowo;
sygnał wraca do modułu~A przewodem sterowniczym przez złącze -X5. Przegrzanie \emph{którejkolwiek}
sekcji (np. zatrzymany wentylator, zasłonięty wylot) $\to$ przerwa $\to$ moc odcięta sprzętowo,
bez udziału procesora.
\item \textbf{-S1 --- krańcówka pokrywy} strefy mocy modułu~A (NC z napędem rolkowym): otwarcie
pokrywy nad częścią 230/400~V przerywa łańcuch.
\item \textbf{-K2 --- potwierdzenie pracy wentylatorów} --- zestyk przekaźnika \emph{prądowego}
(cewka pomiarowa I$>$ w szereg z zasilaniem -M3/-M4). Styk zamyka się tylko, gdy wentylatory
\emph{fizycznie pobierają prąd} --- czyli realnie się kręcą, a nie tylko ,,dostały rozkaz''.
Wybór przekaźnika prądowego zamiast zwykłego stycznika pomocniczego eliminuje fałszywe
potwierdzenie przy przepalonym silniku lub urwanym przewodzie.
\item \textbf{-K10 --- zezwolenie MCU} --- przekaźnik, którego cewkę zasila wyjście cyfrowe -A1
(przez optoizolację). Firmware zamyka -K10 dopiero po pozytywnym self-teście i tylko na czas
kroków wymagających mocy; każda anomalia programowa = zdjęcie zezwolenia.
\item \textbf{-A2 --- watchdog impulsowy} --- człon sprzętowy (ładowanie kondensatora ciągiem
impulsów z -A1, typ.\ kilka Hz; kondensator trzyma tranzystor/przekaźnik zamknięty). Gdy firmware
się zawiesi, impulsy znikają, kondensator się rozładowuje, styk -A2 otwiera łańcuch --- nawet
jeśli ,,zamrożone'' wyjście -K10 zostało w stanie wysokim. To zabezpieczenie przed najgroźniejszą
awarią programową: nie crashem, lecz zawieszeniem z zatrzaśniętymi wyjściami.
\item \textbf{-S2 START $\parallel$ -RB:13-14 samopodtrzymanie} --- przycisk zwierny (NO)
równolegle ze stykiem pomocniczym -RB. Po skompletowaniu całego łańcucha operator musi
\emph{fizycznie} nacisnąć START, by wzbudzić -RB; potem -RB trzyma się własnym stykiem.
Po każdym rozpadzie (E-STOP, termik, watchdog\ldots) łańcuch \textbf{nie wraca sam} --- wymaga
świadomej akcji człowieka. To eliminuje samoczynny restart mocy np. po ostygnięciu termika.
\item \textbf{-RB --- przekaźnik bezpieczeństwa} ze stykami \emph{prowadzonymi wymuszenie}
(mechanically linked): styki NO i NC są sprzężone mechanicznie tak, że nie mogą być jednocześnie
zamknięte; sklejenie styku głównego jest wykrywalne po styku lustrzanym. Klasa wyżej niż zwykły
przekaźnik; tańsza alternatywa dla pełnego modułu PNOZ (do decyzji przy doborze finalnym).
\end{enumerate}

\subsection{Szczebel 2: styk -RB $\to$ cewka -K1}
Osobny szczebel: +24~V $\to$ zestyk -RB:23-24 (wymuszony) $\to$ cewka -K1 $\to$ 0~V.
Rozdzielenie ,,łańcuch zbiera zezwolenia'' / ,,styk mocy steruje stycznikiem'' pozwala
w przyszłości dołożyć redundancję (drugi tor, kontrola sklejenia) bez przebudowy.

\subsection{Diagnostyka łańcucha przez MCU}
Stan członów -S0, -F11.x, -S1, -K2 jest czytany przez -A1 \emph{przez optoizolację} wyłącznie
informacyjnie (komunikat na HMI ,,otwarta pokrywa'', ,,termik sekcji L2''\ldots). Wejścia odczytowe
nie ingerują w tor prądowy łańcucha --- podsłuchują napięcie między członami. Uzupełniająco,
\emph{poza} łańcuchem, działają: -B1 jako RCM (przez -K10), bezpieczniki gG, rozłącznik -Q0
i blokada wtyku -Y1.

% =====================================================================
\section{Arkusz E3 --- zasilanie, sterowanie, moduł obciążenia}
\label{sec:e3}

\subsection{-X4 / -G1 --- zasilanie logiki}
Własny przewód z wtykiem Schuko (-X4) do gniazda instalacji zakładu; zasilacz -G1 daje
\textbf{24~V} (łańcuch bezpieczeństwa, cewki przekaźników/styczników, -Y1) oraz \textbf{5~V
i 3,3~V} (ESP32, czujniki, HMI, serwa). \textbf{Logika nigdy nie jest zasilana z badanej
ładowarki} --- urządzenie musi żyć i nadzorować także wtedy, gdy DUT jest odłączony lub właśnie
zadziałało jego zabezpieczenie. Zanik 24~V = rozpad łańcucha = stan bezpieczny (spójność
z zasadą fail-safe). \popr{-F4, -F5} Dodano zabezpieczenia: -F4 (T2A) przed -G1
i -F5 (T2A) toru wentylatorów --- przewód Schuko chroni instalacja budynku (B16), ale zwarcie
wewnętrzne w zasilaczu/wentylatorze nie powinno czekać na wyłącznik obiektu.

\subsection{Moduł B --- bank obciążenia (K11--K22, E1--E12)}
\label{sec:loadbank}
\textbf{Funkcja:} kontrolowany, rezystancyjny pobór prądu --- urządzenie ,,konsumuje'' energię
jak samochód, zamieniając ją na ciepło w grzałkach żebrowanych 230~V umieszczonych w kanale
z wymuszonym przepływem. \textbf{Struktura:} na każdą fazę cztery stopnie
\textbf{1,0 / 1,5 / 2,0 / 3,0~kW} załączane stycznikami modułowymi 25~A:
L1: -K11..-K14/-E1..-E4, L2: -K15..-K18/-E5..-E8, L3: -K19..-K22/-E9..-E12.
Prądy stopni przy 230~V: 4,3 / 6,5 / 8,7 / 13,0~A; kombinacje dają ok.\ 4,3--32,6~A z krokiem
$\le$2,2~A --- pokrywają typowe progi EVSE 6/10/13/16/20/25/32~A. Uwaga praktyczna: komplet
czterech stopni to nominalnie $\sim$32,6~A ($>$32~A); sekwencer i tak dobiera kombinację
$\le$ deklaracji CP i pilnuje prądu zmierzonego (-T/-U2), a rzeczywisty prąd zależy od napięcia
sieci i tolerancji grzałek --- dobór finalny mocy grzałek doprecyzuje etap zamówień.
\textbf{Dlaczego styczniki, nie SSR:} separacja galwaniczna, odporność na przepięcia, zero
radiatorów; rozdzielczość skokowa wystarcza, a modulacja sekundowa (duty-cycle) też jest możliwa.
\popr{bezpieczniki stopni} Koncepcja (B3) wymagała bezpieczników na stopniach --- na schemacie
ich nie było. W v3 każda gałąź ma wkładkę \textbf{-F12.1..-F12.12} (gG, dobrana do mocy stopnia,
np. 6/10/16~A): zwarcie w pojedynczej grzałce/przewodzie gałęzi przerywa tę gałąź, zamiast czekać
na gG~32~A całej fazy.

\textbf{Wentylatory -M3/-M4} (kanałowe, 2$\times\sim$500~m$^3$/h): przy 11~kW ciągłego grzania
przyrost temperatury powietrza $\sim$33~K --- chłodzenie jest elementem konstrukcyjnym, nie
dodatkiem. \popr{-K3} W v1 zestyk załączający wentylatory był nieopisany, w v2 zniknął zupełnie
(wentylatory pracowałyby zawsze). W v3: \textbf{-K3} (przekaźnik 24~V, sterowany z -A1 przez
ekspander/ULN) załącza wentylatory, a przekaźnik prądowy \textbf{-K2} niezależnie potwierdza
w łańcuchu, że naprawdę pobierają prąd. Sekwencja: firmware włącza -K3 przed testem obciążenia;
bez potwierdzenia -K2 łańcuch nie zamknie -K1.

\textbf{-X5 --- złącze sterownicze modułu B} (wielostykowe przemysłowe, np. Harting/Weipu):
cewki styczników stopni (24~V), pętla termików -F11.x, zasilanie i potwierdzenie wentylatorów.
Moc i sterowanie idą osobnymi przewodami --- rozłączenie CEE nie zostawia ,,żywego'' sterowania.

\subsection{-A3 --- generator prądu upływu (test RDC-DD / RCD)}
\label{sec:a3}
\textbf{Po co:} każda ładowarka AC z RCD typu~A musi mieć układ \textbf{RDC-DD} wykrywający
gładki prąd stały $\ge$6~mA (IEC~62955) --- gładki DC nie jest widziany przez RCD typu~A, a już
kilkanaście mA DC potrafi zablokować (nasycić) jego przekładnik. Test tego zabezpieczenia wymaga
\emph{wymuszenia} kontrolowanego upływu L$\to$PE i zmierzenia, przy jakim prądzie i po jakim
czasie ładowarka się wyłączy. \textbf{Jak działa:} przetwornik DAC (MCP4725, I2C za izolatorem
cyfrowym) zadaje napięcie sterujące; wzmacniacz operacyjny z tranzystorem MOSFET pracuje jako
\textbf{źródło prądowe} 0--20~mA z toru L1 (odczep za -B1, przez zestyk -K30) do PE; prąd
faktyczny mierzy \textbf{bocznik} ze wzmacniaczem izolowanym \textbf{AMC1301}, którego wyjście
czyta ADC w -A1. Procedura: rampa (np. 0,1~mA/s w okolicy progu) $\to$ rejestracja prądu
zadziałania ($\le$6~mA = PASS) i czasu odcięcia (znacznik: zanik faz na optotransoptorach,
rozdz.~\ref{sec:di}). Tryb AC/pulsujący pozwala analogicznie testować RCD typu~A (15/30/150~mA);
pomiar protokołowy RCD w etapie~1 pozostaje przy wzorcowanym UT595. \textbf{Bezpieczeństwo toru:}
szeregowy rezystor mocy ogranicza prąd sprzętowo ($\ll$ progu fibrylacji), a -K30 odłącza
generator zawsze poza testem.

\subsection{-K30..-K33 --- przekaźniki separujące (tryb izolacji)}
\label{sec:k30}
\textbf{Po co:} pomiar rezystancji izolacji (UT595, 250/500~V~DC między zwartymi żyłami czynnymi
a PE) podałby wysokie napięcie na wszystko, co wisi między L/N a PE wewnątrz urządzenia ---
zafałszowałby wynik (równoległe rezystancje) i mógł uszkodzić elektronikę. \textbf{Działanie:}
przed pomiarem izolacji firmware otwiera przekaźniki separujące, odłączając galwanicznie własne
bloki od magistrali. Przyporządkowanie (doprecyzowane w v3): \textbf{-K30} --- tor L1 generatora
-A3; \textbf{-K31/-K32/-K33} --- dzielniki napięciowe L1/L2/L3 układu -U2. Uzasadnienie zakresu:
przy próbie L+N$\,\to\,$PE elementy włączone \emph{między} przewodami czynnymi (dzielniki L--N)
nie widzą napięcia próby, groźne są tylko ścieżki \emph{do PE}; jedyną taką ścieżką roboczą jest
tor -A3, a otwarcie dzielników fazowych dodatkowo zeruje wpływ próby na front-end pomiarowy.
Tor CP/PP nie łączy się z żyłami czynnymi, więc separacji nie wymaga.

\subsection{-Y1 --- blokada elektromagnesowa wtyku Type 2}
\label{sec:y1}
Elektromagnes/napęd rygla w gnieździe -X1, zasilany 24~V, sterowany wyjściem -A1 (opto).
Zamykany po wejściu w stan~C (napięcie na magistrali), otwierany dopiero po potwierdzeniu
beznapięciowości. Uniemożliwia rozłączenie złącza pod obciążeniem (łuk, uszkodzenie styków).

% =====================================================================
\section{Mikrokontroler -A1 i wszystkie urządzenia do niego podpięte}
\label{sec:mcu}

\subsection{-A1: ESP32-S3 --- rola i granice odpowiedzialności}
Płyta logiki z modułem \textbf{ESP32-S3-WROOM} (2 rdzenie LX7 240~MHz, 45~GPIO, WiFi 2,4~GHz,
USB~OTG, 12-bit ADC). Zadania: \textbf{sekwencer} pomiarowy (stany, bramki, blokady kolejności),
\textbf{akwizycja} (CP/PP, metering, temperatury, prąd różnicowy, bocznik -A3),
\textbf{sterowanie} (serwa, stopnie obciążenia, separacja, blokada wtyku, generator),
\textbf{UI} (HMI DWIN, LED, buzzer), \textbf{rejestracja} (microSD, RTC) i \textbf{komunikacja}
(WiFi/MQTT $\to$ Home Assistant). Granica: -A1 \emph{nie jest} barierą bezpieczeństwa --- może
mocy zabronić (zdjąć -K10, przestać karmić -A2), ale nie może jej wymusić wbrew łańcuchowi E2.

Zasilanie: 3,3~V (rdzeń, czujniki) i 5~V (serwa, HMI, przekaźniki wejściowe opto) z -G1;
masa logiki odseparowana od obwodów 24~V (optoizolacja DI/DO) i od potencjałów sieciowych
(izolatory cyfrowe, AMC1301, przekładniki, dzielniki za -K31..33).

\subsection{Tabela zbiorcza urządzeń przy -A1}
\label{sec:tabela-mcu}

\begin{longtable}{@{}p{0.16\textwidth}p{0.13\textwidth}p{0.20\textwidth}p{0.43\textwidth}@{}}
\toprule
\textbf{Aparat} & \textbf{Interfejs} & \textbf{Kierunek} & \textbf{Funkcja w systemie}\\
\midrule
\endhead
-B1 CDSR 0.07-TP & SPI (CS1) & czujnik $\to$ MCU & prąd różnicowy AC/DC magistrali (RCM); weryfikacja -A3\\
-U2 ATM90E36A & SPI (CS2) & czujnik $\to$ MCU & U/I/P/Q/cos$\varphi$/f per faza (z -T1..3 i dzielników); soak\\
karta microSD & SPI (CS3) & MCU $\leftrightarrow$ & protokoły i logi CSV/JSON\\
-B2 MLX90640 & I2C (0x33) & czujnik $\to$ MCU & matryca IR 32$\times$24: $\Delta$T wtyku Type~2, hot-spoty\\
-B3, -B4 MLX90614 & I2C (0x5A/0x5B) & czujnik $\to$ MCU & pirometry punktowe: korpus wtyku, dławik\\
-U3, -U4 MCP23017 & I2C (0x20/0x21) & MCU $\to$ & ekspandery 2$\times$16 wyjść $\to$ -U5/-U6 ULN2803 $\to$ cewki\\
RTC DS3231 & I2C (0x68) & $\leftrightarrow$ & znaczniki czasu protokołów\\
(-A3) MCP4725 & I2C za izolatorem & MCU $\to$ & zadawanie rampy generatora upływu\\
-B5.1..8 DS18B20 & 1-Wire & czujniki $\to$ MCU & temperatury: szyny, -K1, komora grzałek, otoczenie\\
-P1 DWIN 7'' & UART (115200) & $\leftrightarrow$ & HMI: prowadzenie operatora, wpisy wyników, status\\
CP (z -X3) & ADC + capture & sygnał $\to$ MCU & poziom stanu A/B/C/D + duty PWM 1~kHz (deklaracja prądu)\\
PP (z -X3) & ADC & sygnał $\to$ MCU & kod obciążalności kabla (rezystancja)\\
-A3 bocznik/AMC1301 & ADC & sygnał $\to$ MCU & prąd upływu faktycznie płynący\\
-M1, -M2 DS3218 & 2$\times$PWM 50~Hz & MCU $\to$ & serwa pokręteł PM701E (stan / prąd)\\
-K10 & DO (opto) & MCU $\to$ & zezwolenie w łańcuchu E2\\
-A2 watchdog & DO (opto, impulsy) & MCU $\to$ & dowód życia firmware; zanik = rozpad łańcucha\\
-Y1 & DO (opto) & MCU $\to$ & blokada wtyku Type~2\\
-K3 (v3) & przez -U3/-U5 (ULN) & MCU $\to$ & załączenie wentylatorów modułu B\\
-K11..-K22 & przez -U3/4+ULN & MCU $\to$ & stopnie obciążenia L1/L2/L3\\
-K30..-K33 & przez -U3/4+ULN & MCU $\to$ & separacja bloków własnych (tryb izolacji)\\
-S0/-F11.x/-S1/-K2 & DI (opto) & pole $\to$ MCU & diagnostyka członów łańcucha (tylko odczyt)\\
opto faz L1/L2/L3 & 3$\times$DI (opto) & pole $\to$ MCU & obecność napięć; pomiar czasu zadziałania zabezpieczeń\\
LED R/Y/G, buzzer & GPIO / ekspander & MCU $\to$ & sygnalizacja stanu (napięcie / test / gotowość)\\
WiFi / MQTT & radio & $\leftrightarrow$ & podgląd i archiwizacja w Home Assistant\\
\bottomrule
\end{longtable}

\subsection{Magistrala SPI}
\label{sec:mcu-spi}

\textbf{-B1 --- LEM CDSR 0.07-TP (czujnik różnicowy).} Opis zasady --- rozdz.~\ref{sec:b1}.
Podłączenie: wspólne linie SCLK/MOSI/MISO, osobny chip-select; zasilanie 3,3~V wprost z płyty
logiki (czujnik jest konstrukcyjnie izolowany od obwodu pierwotnego --- przewody mocy tylko
przechodzą przez okno). Firmware odpytuje czujnik ciągle (watchdog pomiarowy): wartość AC i DC
osobno; progi alarmowe zależne od kroku sekwencji (podczas testu -A3 --- maska + porównanie
z zadanym prądem).

\textbf{-U2 --- ATM90E36A (analogowy front-end pomiaru energii 3F).}
\label{sec:u2}
Dedykowany układ klasy licznikowej: trzy kanały napięciowe (z dzielników rezystancyjnych
za -K31..33) i trzy prądowe (z przekładników -T1..3). Sam liczy w sprzęcie wartości skuteczne
U/I, moce P/Q/S, cos$\varphi$, częstotliwość i THD per faza --- MCU tylko odczytuje rejestry
po SPI, bez obciążania rdzenia próbkowaniem 50~Hz. Rola w sekwencji: weryfikacja napięć
i kolejności faz po wejściu w stan~C, pomiar prądu każdej fazy w teście obciążenia (spójność
z deklaracją CP), wykrycie spadków napięcia i asymetrii podczas soak, rejestracja przebiegu.

\textbf{Karta microSD} --- na tej samej magistrali (osobny CS): zapis surowych logów i wyników
(CSV/JSON) ze znacznikami RTC; nośnik przeżywa restart i pozwala wygenerować protokół PDF na PC.

\subsection{Magistrala I2C}
\label{sec:mcu-i2c}

\textbf{-B2 --- MLX90640 (matryca termowizyjna 32$\times$24).} Zamontowana nad zatoką gniazda
pojazdowego, polem widzenia ($\sim$55$^\circ$) obejmuje wtyk Type~2 DUT i odcinek kabla.
Termopara-matryca: każdy z 768 pikseli zwraca temperaturę bezdotykowo; odczyt ramek po I2C
(kilka Hz wystarcza --- zjawiska cieplne są wolne). Funkcja: wykrywanie \textbf{hot-spotów}
złącza (klasa usterki U-001 z katalogu diagnostyki) i pomiar $\Delta$T wtyku względem otoczenia
w \textbf{teście soak} (15--30~min pod prądem); kryteria: graniczne $\Delta$T oraz anomalne
tempo narastania d$T$/d$t$ (progi do kalibracji kamerą UTi120S na etapie prototypu).

\textbf{-B3, -B4 --- MLX90614 (pirometry punktowe).} Uzupełniają matrycę tam, gdzie nie sięga jej
pole widzenia (korpus wtyku od drugiej strony, dławik kablowy). Uwaga integracyjna: fabrycznie
oba mają adres 0x5A --- jednemu trzeba przeprogramować adres (EEPROM) albo rozdzielić je na
osobne szyny; na schemacie przyjęto 0x5A/0x5B.

\textbf{-U3, -U4 --- MCP23017 (ekspandery I/O).} Dają 2$\times$16 dodatkowych wyjść, których
ESP32 fizycznie nie ma: 12 cewek stopni (-K11..22), 4 separujące (-K30..33), -K3, rezerwy.
Wyjścia ekspanderów nie ciągną cewek bezpośrednio --- sterują wejściami \textbf{-U5/-U6 ULN2803}
(bufory Darlingtona 500~mA/kanał z diodami gaszącymi), a dopiero te przełączają cewki 24~V
(strona dolna, low-side). Między logiką 3,3~V a stroną 24~V jest optoizolacja --- przepięcia
łączeniowe cewek nie wracają do procesora. Adresy 0x20/0x21 (piny A0--A2).

\textbf{RTC DS3231} --- zegar czasu rzeczywistego z kompensacją temperaturową (TCXO): znaczniki
czasu w protokołach i logach niezależne od WiFi.

\textbf{(-A3) MCP4725} --- 12-bitowy DAC generatora upływu; wisi na I2C \emph{za izolatorem
cyfrowym} (osobna wyspa potencjałowa generatora, rozdz.~\ref{sec:a3}).

\subsection{Magistrala 1-Wire}
\textbf{-B5.1..-B5.8 --- DS18B20 (czujniki temperatury).} Cyfrowe termometry $\pm$0,5~$^\circ$C
na wspólnej jednej linii danych (każdy ma unikalny 64-bitowy adres fabryczny --- stąd jedna
magistrala obsługuje wszystkie). Rozmieszczenie: szyny/listwy toru mocy, obudowa stycznika -K1,
komora grzałek modułu~B (4~szt., po jednej na sekcję + zapas), \textbf{temperatura otoczenia}
(czujnik odniesienia --- bez niego $\Delta$T wtyku nie ma sensu). Funkcja: trend nagrzewania
własnych torów, kryteria pass/fail soak, dane wejściowe progów termicznych firmware (niezależnych
od sprzętowych KSD301, które chronią ,,na twardo'').

\subsection{UART --- HMI}
\textbf{-P1 --- panel DWIN 7'' (DGUS).} Wyświetlacz z własnym kontrolerem grafiki: MCU nie rysuje
pikseli, tylko wymienia po UART (115200) wartości zmiennych i odbiera zdarzenia dotyku ---
ekrany projektuje się narzędziem DGUS (zakład ma je opanowane z wątku ,,chińskie zabawki'',
folder \texttt{chinskie\_zabawki/DWIN\_SET}). HMI: prowadzi operatora krok po kroku (gdzie
przyłączyć UT595, co zmierzyć), przyjmuje ręczne wpisy wyników z przyrządów wzorcowanych
(UT595/UT522 nie mają interfejsów), pokazuje status/awarie. \textbf{Ważne:} fizyczny przycisk
START (-S2) i E-STOP (-S0) są \emph{poza} HMI, w obwodzie E2 --- dotyk ekranu nigdy nie załącza
mocy. Kolumna LED (czerwony = napięcie na magistrali, żółty = test w toku, zielony = gotowość)
i buzzer sterowane z GPIO/ekspandera.

\subsection{Tor analogowy CP/PP (ADC + wejście przechwytujące)}
\label{sec:cp}
Sygnał CP (z gniazd SIGNAL~OUTPUT PM701E) ma poziomy $\pm$12~V --- poza zakresem ADC 0--3,3~V.
Kondycjonowanie: \textbf{dzielnik rezystancyjny} z przesunięciem (mapuje $-$12\ldots+12~V na
0\ldots3,3~V), \textbf{bufor} operacyjny rail-to-rail (niska impedancja dla ADC, brak obciążania
linii CP) i \textbf{ograniczniki diodowe/Zener} (ochrona wejścia przy stanach przejściowych).
Dwa równoległe odbiory:
\begin{itemize}
\item \textbf{ADC}: pomiar poziomu szczytu dodatniego $\to$ rozpoznanie stanu złącza
(+12 = A, +9 = B, +6 = C, +3 = D, 0/$-$12 = E/F);
\item \textbf{komparator $\to$ wejście przechwytujące} (jednostka MCPWM-capture/PCNT):
pomiar wypełnienia PWM 1~kHz z rozdzielczością timera sprzętowego --- z wypełnienia wynika prąd
deklarowany ładowarce: $I\,[\mathrm{A}] = \text{duty}\,[\%]\times 0{,}6$ (np.\ 53\% $\approx$ 32~A,
26,7\% $\approx$ 16~A).
\end{itemize}
Zastosowania: \textbf{weryfikacja pozycji serw} (po komendzie ,,stan C'' firmware czeka na
+6~V na CP --- niezależne potwierdzenie, że pokrętło naprawdę się przestawiło), zapis parametrów
CP do protokołu, wykrywanie błędów zgłaszanych przez ładowarkę. \textbf{PP}: drugi kanał ADC
czyta rezystancję kodującą obciążalność kabla (1,5~k\ohm{}=13~A, 680~\ohm{}=20~A,
220~\ohm{}=32~A, 100~\ohm{}=63~A wg IEC~61851-1). Mikrooscyloskop pozostaje niezależnym
podglądem ręcznym tego samego sygnału.

\subsection{PWM --- serwomechanizmy pokręteł PM701E}
\label{sec:serwa}
\textbf{-M1 (pokrętło stanu BB/C/A/B/D), -M2 (pokrętło prądu NC/13/20/32/63~A)} --- serwa
modelarskie klasy \textbf{DS3218} (20~kg$\cdot$cm, metalowa przekładnia; moment pokręteł
przełączników to typowo 1--3~kg$\cdot$cm $\Rightarrow$ zapas $\ge$5$\times$). Sterowanie:
standardowy sygnał serw 50~Hz, impuls 500--2500~µs $\to$ kąt; zasilanie 5~V z wydzielonej
gałęzi -G1 (szczytowo $\sim$1~A/serwo --- nie z szyny czujników). Montaż nieinwazyjny: rama nad
panelem PM701E, sprzęg przez drukowany adapter i \textbf{sprzęgło Oldham} (kompensuje
niewspółosiowość, chroni oś przełącznika przed siłami promieniowymi). Firmware trzyma tabelę
kalibracji \emph{kąt $\leftrightarrow$ pozycja} (5+5 pozycji dyskretnych) i \textbf{nigdy nie ufa
samemu serwu}: każdą zmianę potwierdza odczytem CP (rozdz.~\ref{sec:cp}); brak potwierdzenia
$\to$ przerwanie sekwencji i zdjęcie zezwolenia. PM701E pozostaje nienaruszony (gwarancja,
kalibracja).

\subsection{Wyjścia cyfrowe (DO, przez optoizolację)}
\label{sec:do}
Trzy bezpośrednie, rozdzielone wyjścia (każde ze swoim transoptorem, klucz low-side na 24~V):
\begin{itemize}
\item \textbf{DO1 $\to$ -K10}: zezwolenie --- styk w łańcuchu E2;
\item \textbf{DO2 $\to$ -A2}: ciąg impulsów watchdog (oddzielny od DO1! zamrożenie DO1
w stanie ,,1'' nie utrzyma mocy, bo brak impulsów rozłączy -A2);
\item \textbf{DO3 $\to$ -Y1}: blokada wtyku.
\end{itemize}
\popr{rozdzielenie DO} W v2 linie DO do -K10, -A2 i -Y1 narysowano jako jeden wspólny węzeł ---
elektrycznie oznaczałoby to jedno wyjście sterujące naraz zezwoleniem, watchdogiem i blokadą
(bez sensu funkcjonalnie: watchdog wymaga impulsów, zezwolenie stanu ciągłego). W v3 ---
trzy niezależne tory zgodnie z v1 i koncepcją.
Pozostałe wyjścia (cewki -K11..22, -K30..33, -K3) idą ścieżką
\textbf{I2C $\to$ -U3/-U4 $\to$ -U5/-U6 $\to$ cewki 24~V} (rozdz.~\ref{sec:mcu-i2c}).

\subsection{Wejścia cyfrowe (DI, przez optoizolację)}
\label{sec:di}
\begin{itemize}
\item \textbf{Diagnostyka łańcucha}: stany -S0, pętli -F11.x, -S1, -K2 --- wyłącznie odczyt
(HMI pokazuje, \emph{który} człon rozłączył); tor łańcucha pozostaje nietknięty.
\item \textbf{Optotransoptory obecności faz L1/L2/L3} (wejściowe AC, typu H11AA1, za rezystorami):
przy obecności napięcia dają ciąg impulsów 100~Hz; zanik impulsów $=$ zanik fazy. Kluczowe dla
pomiaru \textbf{czasu zadziałania} zabezpieczeń DUT: przy teście RDC-DD/RCD i testach błędów
CP/PE firmware stempluje czas od bodźca do zaniku napięcia (rozdzielczość pojedynczych ms ---
wystarczająca wobec limitów normowych rzędu dziesiątek/setek ms).
\end{itemize}

\subsection{Pozostałe zasoby -A1}
\textbf{WiFi/MQTT} --- telemetria i archiwizacja w zakładowym Home Assistant (spójnie z wątkiem
integracji Tuya/HA); brak sieci nie blokuje pomiarów (SD jest źródłem prawdy).
\textbf{USB~OTG} --- serwis/aktualizacja firmware. Orientacyjny budżet wyprowadzeń: SPI(4)+3$\times$CS,
I2C(2), UART HMI(2), 1-Wire(1), ADC(3: CP, PP, bocznik), capture(1), PWM(2), DO(3), DI($\sim$7),
LED/buzzer(4), SD-detect itp. --- ok.\ 33--36 z 45 dostępnych GPIO; rezerwa pozostaje na aktuatory
przycisków PM701E (etap rozszerzenia).

% =====================================================================
\section{Jak bloki grają razem --- sekwencja w ujęciu elektrycznym}
\label{sec:sekwencja}

\begin{enumerate}
\item \textbf{Self-test}: -A1 czyta łańcuch (DI), próbuje -K3/-K2 (wentylatory), odpytuje
-B1/-U2/-B2..5/-P1, jeździ serwami do pozycji spoczynkowych (BB/NC) z potwierdzeniem CP.
\item \textbf{Tryb izolacji}: -K30..33 otwarte, moduł~B odłączony (-K1 otwarty) --- UT595 na
gniazdach bananowych PM701E mierzy izolację; wynik wpisany na HMI otwiera bramkę.
\item \textbf{Ciągłość PE}: UT595 (wyjście PE) --- wpis, bramka.
\item \textbf{Stan B $\to$ C}: -M1 przestawia pokrętło, CP potwierdza; w stanie C ładowarka
zamyka swój stycznik --- -U2 weryfikuje napięcia i kolejność faz, -Y1 rygluje wtyk.
\item \textbf{Zabezpieczenia}: UT595 protokołowo (RCD~A); -A3 automatycznie (rampa 6~mA DC),
czas zadziałania z optotransoptorów faz; po każdym zadziałaniu prowadzony reset DUT.
\item \textbf{Pętla Zs}: UT595 --- wpis.
\item \textbf{Obciążenie/soak}: -M2 ustawia deklarację, łańcuch zamknięty + START $\to$ -K1;
stopnie 25/50/75/100\% deklaracji (kombinacje $\le$ CP); -U2 loguje U/I/P, -B2/-B5 temperatury;
kryteria pass/fail.
\item \textbf{Testy błędów}: przyciski CP~Error/PE~Error na PM701E (etap 1 ręcznie) --- czas
reakcji DUT mierzony jak wyżej.
\item \textbf{Uziom}: UT522 na zaciskach E/S/H (poza torem Type~2).
\item \textbf{Raport}: SD + MQTT; protokół PDF na PC.
\end{enumerate}
Warunki ciągłe (każdy krok): CP zgodne z komendą, RCM poniżej progu, termiki zwarte, wentylatory
potwierdzone, E-STOP zwolniony --- naruszenie = przerwanie i stan bezpieczny.

% =====================================================================
\section{Wykryte problemy i poprawki wprowadzone w wersji 3}
\label{sec:poprawki}

Przegląd v1/v2 wykazał sześć problemów wymagających korekty oraz trzy braki o charakterze
uzupełnień bezpieczeństwa. Wszystkie wprowadzono w \textbf{schemacie elektrycznym v3}
(\texttt{AMPERE\_POINT\_custom\_device\_schemat\_elektryczny\_v3\_EPLAN}); v1/v2 pozostają
w folderze bez zmian (zasada wersjonowania).

\begin{longtable}{@{}p{0.04\textwidth}p{0.06\textwidth}p{0.42\textwidth}p{0.40\textwidth}@{}}
\toprule
\textbf{\#} & \textbf{Ark.} & \textbf{Problem (v1/v2)} & \textbf{Poprawka (v3)}\\
\midrule
\endhead
1 & E1 & Odczep -X3 tylko 5-żyłowy (L1, N, PE, CP, PP) --- gniazda L2/L3 PM701E martwe;
niespójność z koncepcją (7 żył) i z pomiarami 3F & Odczep 7-żyłowy: L1, L2, L3, N, PE, CP, PP;
numeracja zacisków 1--7\\
2 & E3 & Wspólny węzeł DO dla -K10, -A2 i -Y1 (jedno wyjście sterujące zezwoleniem, watchdogiem
i blokadą) --- błąd rysunkowy v2 sprzeczny z v1 i koncepcją & Rozdzielone wyjścia DO1..DO4
(osobne transoptory i tory)\\
3 & E3 & Brak aparatu załączającego wentylatory -M3/-M4 (v2: pracowałyby zawsze przy zasilonym
-X4; v1: zestyk nieopisany) & Nowy przekaźnik \textbf{-K3} (cewka 24~V, sterowanie z -A1);
-K2 pozostaje wyłącznie potwierdzeniem prądowym w łańcuchu\\
4 & E3 & Brak bezpieczników gałęziowych stopni obciążenia --- wymaganych w koncepcji (B3);
zwarcie w gałęzi czekałoby na gG~32~A całej fazy & Wkładki \textbf{-F12.1..-F12.12} w każdej
gałęzi stopnia (dobór do mocy stopnia)\\
5 & E1/E3 & Punkt przyłączenia generatora -A3 wskazany przy -F1 (\texttt{/E1.2}), czyli
\emph{przed} oknem -B1 --- CDSR nie widziałby wymuszanego upływu; utrata weryfikacji krzyżowej
i nadzoru RCM & Odczep -A3 przeniesiony \textbf{za okno -B1} (między -B1 a -T1), zaznaczony
na E1\\
6 & E2 & Odsyłacz styku samopodtrzymania -RB:13-14 wskazywał kolumnę 3 zamiast 4 (kosmetyka
nawigacji) & Odsyłacze uzgodnione z siatką (13-14 \texttt{/E2.4})\\
7 & E1/E2 & \emph{Luka bezpieczeństwa:} E-STOP otwierał tylko -K1; magistrala i gniazda bananowe
zostawały pod napięciem od DUT w stanie C & Nowy przekaźnik \textbf{-K4}: zestyk w linii CP
odczepu, cewka przez -S0:21-22; E-STOP/zanik 24~V $\to$ przerwanie CP $\to$ DUT sam otwiera
stycznik (mechanizm ,,CP Error''); mostek serwisowy na -X3\\
8 & E3 & Brak zabezpieczeń wewnętrznych toru logiki (zasilacz, wentylatory na wspólnym
przewodzie Schuko) & Wkładki \textbf{-F4} (T2A, przed -G1) i \textbf{-F5} (T2A, tor
wentylatorów)\\
9 & E1 & Brak ostrzeżenia, że -Q0 nie odłącza odczepu -X3/PM701E (gniazda bananowe pod napięciem
w stanie C) & Adnotacja na arkuszu + wymóg tabliczki ostrzegawczej na panelu; systemowo
domyka to -K4 (poz.~7)\\
\bottomrule
\end{longtable}

Doprecyzowania bez zmiany połączeń: przyporządkowanie -K30..33 (rozdz.~\ref{sec:k30});
uwaga o adresach MLX90614 (rozdz.~\ref{sec:mcu-i2c}); uwaga o $\sim$32,6~A przy komplecie stopni
(rozdz.~\ref{sec:loadbank}).

\section{Uwagi końcowe i punkty otwarte}
\begin{itemize}
\item Wartości i typy aparatów pozostają \textbf{robocze do etapu~0} (prototyp stołowy):
zwłaszcza test serw na fizycznym PM701E, próbka CDSR/RCM14, prototyp jednej sekcji grzałkowej.
\item Nadal brakuje instrukcji PM701E (funkcje Touch/PE~Pre-Test, prąd znamionowy gniazd
bananowych --- istotny dla Zs dużym prądem) --- otwarty punkt z koncepcji.
\item Wersja formalna schematu powstanie w EPLAN (projekt \texttt{AP-CD-001}); arkusze v3
odwzorowują docelowy układ stron, co uprości przeniesienie.
\item Pomiary protokołowe wykonują przyrządy wzorcowane (UT595/UT522) --- bloki własne
uzupełniają je tam, gdzie przyrządy gotowe nie sięgają (RDC-DD 6~mA~DC, obciążenie, termika,
automatyzacja) --- zgodnie z zatwierdzonym punktem Z4.
\end{itemize}

\end{document}
